\begin{textsample}{POS Dim 6 – global\_north\_2024\_04 – Score 9.00 – global\_north\_2024\_04/107333972.txt}  \label{ex:f6_pos_077}
Cost of US effort to build humanitarian \textbf{aid} \textbf{pier} off Gaza expected to top \$180M The Trident \textbf{pier} rests on the shore of Fort Story, Va., during the preliminary stages of the Joint Logistics-Over-the-Shore exercise, Aug.17, 2012.
DVIDS President Joe Biden ' s plan to build a humanitarian \textbf{pier} floating off the \textbf{coast} of Gaza that could enable delivery of food, water and medicine into the devastated region is expect to cost at least \$180 million and could top \$200 million, ABC News has learned.
The price tag was described by two people familiar with the initial estimate, which has not been released by U.S.
Central Command.
The tally is expected to fluctuate as U.S. officials scramble to finalize key details on the project, including which humanitarian relief organizations and foreign governments are willing to help carry the \textbf{shipments} to shore and distribute them to ease the humanitarian crisis unfolding in Gaza.
Officials also continue to discuss how to protect service members who will be operating three miles offshore of Gaza, where Hamas is believed to still operate. @ @ @ @ @ @ @ @ @ @ and Navy \textbf{ships} and will involve some 1,000 U.S. military troops -- is on track to become operational in early May, enabling the delivery of some 2 million \textbf{meals} a day.
Trucks transporting heavy equipment and large cement blocks for use in constructing a new \textbf{pier} on the Gaza City seashore, on March 15, 2024."No U.S. boots will be on the ground,"Biden promised when announcing the project in his State of the Union speech last month."A temporary \textbf{pier} will enable a massive increase in the amount of humanitarian assistance getting into Gaza every day."The effort has largely been seen as a political move by the president, who faced criticism for not doing more to try to rein in Israel ' s destruction of Gaza following the Oct. 7 Hamas attacks inside Israel and for not forcing Israel to allow more humanitarian \textbf{aid} to get through.
White House officials have acknowledged the \textbf{pier} is still weeks away and not as efficient in delivering \textbf{aid} as via ground convoys.
But they @ @ @ @ @ @ @ @ @ @ open up every \textbf{aid} route possible to address potential \textbf{famine} in Gaza.
For its part, Israel defended its initial refusal to open more humanitarian channels in recent months, noting a need to screen supplies carefully to ensure they don't help Hamas.
The floating dock is expected to be nearly the size of a football field -- about 97 feet wide and 270 feet long -- stationed about three miles offshore.
Container \textbf{ships} would screen their \textbf{cargo} in Cyprus before taking it to the floating dock and \textbf{unloading} it.
From there, the \textbf{aid} would be moved aboard small Army ferries that would transport it to an 1,800-foot"trident"\textbf{pier} that connects to shore.
But with U.S. troops not allowed to go onshore, it ' s still not clear who will bring the \textbf{cargo} from the \textbf{pier} to shore and then distribute it.
Officials have said only that it ' s working with"regional partners"on a solution that will ensure no American boots are on the ground in Gaza.
Deputy Defense Secretary Sabrina @ @ @ @ @ @ @ @ @ @ working"around the clock"to find partners and set up the system."Still no boots on the ground.
That is the policy that has been set by the president.
We will not have boots on the ground when it comes to setting up this \textbf{pier},"she said.

% matched lemmas: aid, cargo, coast, famine, meal, pier, ship, shipment, unload
\end{textsample}
